\documentclass[11pt,a4paper]{article}
\usepackage[T1]{fontenc}
\usepackage[utf8]{inputenc}
\usepackage{lmodern,microtype}
\usepackage[a4paper,textwidth=151mm,top=23mm,bottom=25mm]{geometry}
\usepackage{amsmath,amssymb,amsthm,mathtools}
\usepackage{array,booktabs,needspace,float}
\usepackage[font=small,labelfont=bf]{caption}
\usepackage[section]{placeins}
\usepackage{tikz}
\usetikzlibrary{arrows.meta,positioning,calc,backgrounds,fit,shapes.geometric}
\usepackage[most]{tcolorbox}
\usepackage{titlesec}
\usepackage[colorlinks=true]{hyperref}
\definecolor{articleblue}{HTML}{234D72}
\hypersetup{linkcolor=articleblue,citecolor=articleblue,urlcolor=articleblue,
 pdftitle={Piggyback rules for quantum--classical gaps}}
\urlstyle{same}
\titleformat{\section}{\Large\sffamily\bfseries}{\thesection}{.65em}{}
\titleformat{\subsection}{\large\sffamily\bfseries}{\thesubsection}{.65em}{}
\newtcolorbox{piggyback}[1]{colback=blue!3!white,colframe=articleblue,
 colbacktitle=articleblue,fonttitle=\sffamily\bfseries,title={#1},
 boxrule=.55pt,arc=1mm,left=2.3mm,right=2.3mm,top=1.3mm,bottom=1.3mm,
 before skip=8pt,after skip=8pt}
\newtheorem{proposition}{Proposition}
\newcommand{\Cset}{\mathcal C}
\newcommand{\Qset}{\mathcal Q}
\DeclareMathOperator{\gap}{\text{\normalfont\sffamily\bfseries Gap}}
\newcommand{\Gtarget}{G_{\text{target}}}
\newcommand{\Gfamiliar}{G_{\text{familiar}}}
\newcommand{\Ptarget}{P_{\text{target}}}
\newcommand{\Pfamiliar}{P_{\text{familiar}}}
\newcommand{\ind}{\mathrel{\perp\!\!\!\perp}}
\newcommand{\pa}{\operatorname{pa}}
\newcommand{\ch}{\operatorname{ch}}
\newcommand{\qcn}{QC-gap-non-decreasing}
\setlength{\emergencystretch}{1.4em}
\renewcommand{\floatpagefraction}{.76}
\renewcommand{\topfraction}{.91}
\renewcommand{\bottomfraction}{.82}
\renewcommand{\textfraction}{.08}
\AddToHook{cmd/subsection/before}{\FloatBarrier\Needspace{12\baselineskip}}
% Coordinate unit: mm. Each declaration lists visible vertices, sources, then
% every directed edge, including source-to-recipient edges. Qn denotes psi_n,
% Cn denotes lambda_n, and Bp is the second output field B'. Figure panels
% use these declarations verbatim, without scaling the picture or its text.
\definecolor{censusBrown}{HTML}{886340}
\definecolor{censusAmber}{HTML}{83500D}
\definecolor{censusVisible}{HTML}{FFFBF2}
\definecolor{censusGreen}{HTML}{245E43}
\definecolor{censusGreenLine}{HTML}{42805A}
\definecolor{censusMint}{HTML}{E5F3E7}
\definecolor{censusRed}{HTML}{913E46}
\definecolor{censusRedLine}{HTML}{AF5960}
\definecolor{censusRose}{HTML}{FBE8E7}
\definecolor{censusInk}{HTML}{41454A}
\tikzset{
  census node/.style={inner sep=0pt,outer sep=.23pt,minimum size=8mm,line width=.55pt,
    font=\normalfont\fontsize{10.5}{11}\selectfont},
  census visible/.style={census node,rectangle,rounded corners=.8mm,
    draw=censusBrown,
    fill=censusVisible,text=censusAmber},
  census latent/.style={census node,circle},
  census Q/.style={census latent,draw=censusGreenLine,fill=censusMint,text=censusGreen},
  census C/.style={census latent,draw=censusRedLine,fill=censusRose,text=censusRed},
  census edge/.style={-{Latex[length=1.65mm,width=1.08mm]},line width=.58pt,
    shorten >=.25pt,shorten <=.3pt}
}
\newcommand{\censusLabelBp}{\mathbf B'}
\newcommand{\censusVisibleLabel}[1]{%
  \ifcsname censusLabel#1\endcsname\csname censusLabel#1\endcsname
  \else\mathbf{#1}\fi}
\def\censusSymbolQ{\psi}
\def\censusSymbolC{\lambda}
\def\censusLineQ{censusGreenLine}
\def\censusLineC{censusRedLine}
\def\censusSource#1#2\relax#3#4{%
  \node[census #1] (#1#2) at (#3,#4) {$\csname censusSymbol#1\endcsname_{#2}$};}
\newcommand{\censusPlot}[3]{%
  \begin{tikzpicture}[x=1mm,y=1mm,baseline=(current bounding box.center)]
    \path[use as bounding box] (-3,-4) rectangle (64,48);
    \foreach \v/\xx/\yy in {#1}{
      \node[census visible] (\v) at (\xx,\yy) {$\censusVisibleLabel{\v}$};}
    \foreach \v/\xx/\yy in {#2}{\expandafter\censusSource\v\relax{\xx}{\yy}}
    \begin{scope}[on background layer]
      \foreach \tail/\head in {#3}{
        \def\edgecolor{censusInk}%
        \expandafter\censusEdgeColor\tail\relax
        \draw[census edge,draw=\edgecolor] (\tail) -- (\head);}
    \end{scope}
  \end{tikzpicture}}
% Visible C and classical source Cn share an initial. Only sources have a suffix
% after Q or C. The primed visible Bp always uses the visible-arrow color.
\def\censusEdgeColor#1#2\relax{%
  \if\relax\detokenize{#2}\relax\else
    \ifcsname censusLine#1\endcsname\edef\edgecolor{\csname censusLine#1\endcsname}\fi
  \fi}
\newcommand{\Diagram}[4]{%
  \expandafter\newcommand\csname censusPanel#1\endcsname{\censusPlot{#2}{#3}{#4}}}
\newcommand{\censusPanel}[2]{%
  \begin{minipage}[c]{.46\linewidth}\centering
    \csname censusPanel#1\endcsname\par\vspace{.7mm}
    {\small\sffamily #2\par}
  \end{minipage}}
\newcommand{\censusPair}[6]{%
  \expandafter\newcommand\csname censusFigure#1\endcsname{%
    \begin{figure}[htbp]
      \centering
      \censusPanel{#2}{\textbf{(a)}\enspace #3}\hfill
      \makebox[0pt]{\raisebox{2ex}{$\longrightarrow$}}\hfill
      \censusPanel{#4}{\textbf{(b)}\enspace #5}
      \caption{#6}\label{fig:#1}
    \end{figure}}}
\newcommand{\censusfigure}[1]{\csname censusFigure#1\endcsname}

\Diagram{deg-before}
  {C/30/38,B/8/7,S/52/7}
  {Q1/9/29,Q3/51/29,Q2/30/5}
  {Q1/B,Q1/C,Q3/C,Q3/S,Q2/B,Q2/S}
\Diagram{deg-after}
  {C/30/38,B/8/7,S/52/7}
  {C1/9/29,C3/51/29,Q2/30/5}
  {C1/B,C1/C,C3/C,C3/S,Q2/B,Q2/S}
\censusPair{degradation}{deg-before}{All-quantum triangle}
  {deg-after}{Two sources made classical}
  {Degrading \(\psi_1\) on \(BC\) and \(\psi_3\) on \(CS\) gives the
   classical sources \(\lambda_1,\lambda_3\). Laboratory \(C\) records
   both settings, while \(B,S\) measure \(\psi_2\). This is the final
   step after the upward lookup in Figure~\ref{fig:upward} and fixing \(A\).}

\Diagram{fix-before}
  {A/10/36,B/13/8,C/35/8,D/54/30}
  {Q1/2/21,Q2/32/31}
  {A/B,B/C,C/D,Q1/A,Q1/B,Q2/A,Q2/C,Q2/D}
\Diagram{fix-after}
  {B/13/8,C/35/8,D/54/30}
  {Q2/32/31}
  {B/C,C/D,Q2/C,Q2/D}
\censusPair{fixing}{fix-before}{Two source inputs at \(A\)}{fix-after}{Fix \(A\)}
  {Fixing \(A\) to a constant removes its message to \(B\).
   Source \(\psi_2\) now supplies only \(C,D\), and the surviving share of
   \(\psi_1\) can be prepared locally at \(B\). The setting \(B\), wire
   \(C\to D\), and shared state on \(CD\) give an instrumental seed.}

\Diagram{naive-before}
  {A/8/8,B/52/8,C/30/26,D/30/43}
  {Q1/30/4,Q2/8/28,Q3/52/28}
  {C/D,Q1/A,Q1/B,Q2/A,Q2/C,Q3/B,Q3/C,Q3/D}
\Diagram{naive-after}
  {A/8/8,B/52/8,D/30/43}
  {Q1/30/4,C2/8/28,Q3/52/28}
  {Q1/A,Q1/B,C2/A,C2/D,Q3/B,Q3/D}
\censusPair{naive}{naive-before}{Classical message \(C\to D\)}
  {naive-after}{Marginalize \(C\)}
  {Ordinary marginalization makes a classical register from the \(AC\)
   source available to \(A,D\), drawn as \(\lambda_2\). The original
   \(BCD\) source retains the \(B,D\) shares and can also carry its own
   forwarded classical register. Degrading this last quantum source gives
   the triangle seed with setting recorder \(D\).}

\Diagram{tele-before}
  {A/8/35,B/29/22,C/9/6,D/52/9}
  {Q1/48/35,Q2/2/21}
  {A/B,B/C,B/D,Q1/A,Q1/B,Q1/D,Q2/A,Q2/C}
\Diagram{tele-after}
  {B/29/22,C/9/6,D/52/9}
  {Q1/48/35,Q2/2/21}
  {B/C,B/D,Q1/B,Q1/D,Q2/B,Q2/C}
\censusPair{teleportation}{tele-before}{Entanglement shared by \(A,B\)}
  {tele-after}{Marginalize \(A\)}
  {The source \(\psi_1\) on \(ABD\) can supply an entangled pair to
   \(A,B\) alongside the desired state for \(B,D\). Laboratory \(A\)
   uses the pair to teleport its share of \(\psi_2\) to \(B\) and
   sends the correction along \(A\to B\). The two states on \(BC\)
   and \(BD\) give the Evans seed.}

\Diagram{stitch-before}
  {A/52/38,B/30/28,C/8/15,D/52/15}
  {Q1/16/3,Q2/44/3}
  {A/D,B/C,Q1/B,Q1/C,Q2/B,Q2/D}
\Diagram{stitch-after}
  {B/30/28,C/8/15,D/52/15}
  {Q1/16/3,Q2/44/3}
  {B/C,B/D,Q1/B,Q1/C,Q2/B,Q2/D}
\censusPair{stitching}{stitch-before}{Outcome \(B\) and setting \(A\)}
  {stitch-after}{Stitch \(A\) onto \(B\)}
  {The stitch retains the wire \(B\to C\) and redirects \(A\to D\)
   to \(B\to D\). Keeping records with \(A=B\) supplies the Evans
   experiment on the right, with two independent sources incident to \(B\).}

\Diagram{post-before}
  {B/7/38,C/14/9,D/49/9,A/46/38}
  {Q1/25/26,Q2/57/25}
  {A/B,B/C,C/D,Q1/A,Q1/C,Q2/A,Q2/D}
\Diagram{post-after}
  {B/7/38,C/14/9,D/49/9}
  {Q3/33/28}
  {B/C,C/D,Q3/C,Q3/D}
\censusPair{postselection}{post-before}{Two independent states at \(A\)}
  {post-after}{Postselect at \(A\)}
  {Postselecting a joint-measurement outcome at \(A\) leaves \(C,D\)
   with a shared state, drawn as \(\psi_3\). On the selected branch
   \(B\) receives a fixed value from \(A\). It supplies a setting to \(C\)
   in the resulting instrumental experiment.}

\Diagram{split-before}
  {A/12/23,B/26/42,C/36/23,D/47/23}
  {Q1/0/23,Q2/24/23,Q3/60/23}
  {A/B,B/C,C/D,Q1/A,Q1/B,Q2/A,Q2/C,Q3/B,Q3/D}
\Diagram{split-after}
  {A/12/23,B/26/42,Bp/26/3,C/36/23,D/47/23}
  {Q1/0/23,Q2/24/23,Q3/60/23}
  {A/B,A/Bp,B/C,Bp/C,C/D,Q1/A,Q1/B,Q1/Bp,Q2/A,Q2/C,Q3/B,Q3/Bp,Q3/D}
\censusPair{splitting}{split-before}{One reported field at \(B\)}
  {split-after}{Split \(B\) into two fields}
  {Splitting \(B\) into \(B,B'\) gives both fields the original inputs
   of \(B\) and sends both outputs to \(C\). Either source enlarged by the
   split can supply the shared private randomness. Panel (b) shows the
   intermediate graph before parent prediction restricts \(B'\).}

\Diagram{parent-before}
  {A/12/23,B/26/42,Bp/26/3,C/36/23,D/47/23}
  {Q1/0/23,Q2/24/23,Q3/60/23}
  {A/B,A/Bp,B/C,Bp/C,C/D,Q1/A,Q1/B,Q2/A,Q2/C,Q3/B,Q3/D}
\Diagram{parent-after}
  {B/16/30,C/35/15,D/57/27,Bp/14/0}
  {Q2/35/33,Q3/38/44,C2/14/15}
  {B/C,Bp/C,C/D,Q2/B,Q2/C,Q3/B,Q3/D,C2/B,C2/Bp,C2/C}
\censusPair{parent}{parent-before}{Only \(A\) supplies \(B'\)}
  {parent-after}{Marginalize predictor \(A\)}
  {Laboratory \(A\) predicts \(B'\). Panel (a) shows the intermediate
   graph obtained by removing the two source inputs to \(B'\) from
   Figure~\ref{fig:splitting}(b). Marginalizing \(A\) gives panel (b).
   Laboratory \(A\) uses \(\psi_1\) to teleport the share of \(\psi_2\) to
   \(B\) and forwards a classical register from the \(AC\) source, giving
   \(\lambda_2\) on \(B,B',C\). The next census move marginalizes \(B\)
   to reach an instrumental seed.}

\Diagram{dsep-before}
  {A/8/41,B/40/31,C/52/3,D/20/13}
  {Q1/52/41,Q2/20/31,Q3/8/3,Q4/40/13}
  {Q1/A,Q1/B,Q1/C,Q2/A,Q2/B,Q2/D,Q3/A,Q3/C,Q3/D,Q4/B,Q4/C,Q4/D}
\Diagram{dsep-after}
  {B/30/38,C/8/7,D/52/7}
  {C1/9/29,C2/51/29,Q3/30/5}
  {C1/B,C1/C,C2/B,C2/D,Q3/C,Q3/D}
\censusPair{dsep}{dsep-before}{Four independent triple sources}
  {dsep-after}{Restrict \(B\) and forget \(A\)}
  {Laboratory \(A\) predicts \(B\) using the sources \(\psi_1,\psi_2\).
   The d-separation test allows \(B\) to drop the input from \(\psi_4\)
   and to receive only classical data from \(\psi_1,\psi_2\). Removing
   \(A\) leaves both \(\psi_3\) and \(\psi_4\) on \(CD\). Panel (b) draws
   the tensor product as \(\psi_3\).}

\Diagram{kpc-before}
  {C/9/37,D/9/9,E/50/37,F/50/9}
  {Q1/57/23,Q2/30/5}
  {C/D,C/E,D/E,Q1/E,Q1/F,Q2/D,Q2/F}
\Diagram{kpc-after}
  {C/9/37,D/9/9,E/50/37,F/50/9}
  {C1/57/23,Q2/30/5}
  {C/D,C1/E,C1/F,Q2/D,Q2/F}
\censusPair{semigraphoid}{kpc-before}{\(F\) predicts \(E\)}
  {kpc-after}{Restrict the inputs of \(E\)}
  {When \(F\) predicts \(E\), the semigraphoid certificate allows
   deletion of \(C\to E\) and \(D\to E\). In the Bell graph, \(\lambda_1\)
   supplies the setting at \(E,F\), and \(\psi_2\) supplies the state
   measured at \(D,F\). The proof constructs a fresh source on \(EF\)
   in the classical model.}
\censusPair{entropic}{kpc-before}{Prediction constraint \(H(E\mid F)=0\)}
  {kpc-after}{Bell experiment}
  {The entropy program certifies the deletion of \(C\to E\) and
   \(D\to E\), the same census move as Figure~\ref{fig:semigraphoid}.
   The proof uses \(H(E\mid F)=0\) and the observable independence
   \(I(E;CD)=0\) to justify a fresh classical source on \(EF\). The
   census records the semigraphoid certificate for the move.}

\Diagram{up-before}
  {A/30/28,B/30/7,C/9/38,S/51/38}
  {Q1/12/24,Q2/48/24,Q3/30/42}
  {A/B,Q1/A,Q1/B,Q1/C,Q2/A,Q2/B,Q2/S,Q3/C,Q3/S}
\Diagram{up-after}
  {A/30/28,B/30/7,C/9/38,S/51/38}
  {Q1/12/24,Q2/48/24,Q3/30/42}
  {Q1/A,Q1/B,Q1/C,Q2/A,Q2/B,Q2/S,Q3/C,Q3/S}
\censusPair{upward}{up-before}{Target with \(A\to B\)}
  {up-after}{Remove the visible arrow}
  {The two graphs share every source and have the same classical model set.
   The quantum model set on the right is strictly smaller. An upward
   lookup transfers a gap from right to left. Fixing \(A\) on either side
   gives an all-quantum triangle on \(B,C,S\).}


\begin{document}
\begin{center}
{\LARGE\sffamily\bfseries Piggyback rules for quantum--classical gaps\par}
\medskip
{\large Graph transformations and perfect prediction in causal networks\par}
\end{center}
\begin{abstract}
A causal network has a quantum--classical gap when an experiment with independent
quantum sources produces correlations that independent classical sources cannot
reproduce on the same network. A \emph{piggyback} uses a witness from one graph to
establish a gap on another. The elementary rules use source degradation, value
fixing, marginalization with or without teleportation, stitching, or postselection.
The Fritz rules use outputs that another laboratory can predict exactly. We
illustrate each rule with a step from the four-node census, where piggybacks and
known witnesses establish gaps in 380 of 426 candidate classes.
\end{abstract}

\section{From a familiar witness to a new gap}\label{sec:setting}
In a Bell experiment, two laboratories measure a shared state in response to
settings. The resulting distribution can violate a Bell inequality and so
witness a quantum--classical gap \cite{Bell}. Suppose we want to establish a
gap in a larger network. Can the larger network reproduce the Bell experiment?
Could the extra arrows and sources give the Bell correlations a classical
explanation? A piggyback answers both questions.

The \emph{target graph} \(\Gtarget\) is the graph whose gap we seek to establish.
A piggyback turns the target graph into a \emph{familiar graph} \(\Gfamiliar\).
Write \(\gap(G)\) for the statement
that some quantum experiment on \(G\) produces a distribution with no classical
explanation on \(G\). The graph operation and the gap implication run in opposite
directions:
\begin{equation}
  \begin{aligned}
  \text{Graph operation} &\quad \Gtarget \longrightarrow \Gfamiliar ,\\
  \text{Gap transfer} &\quad \gap(\Gfamiliar) \implies \gap(\Gtarget) .
  \end{aligned}
  \label{eq:direction}
\end{equation}
We call a witness transfer \emph{\qcn} when a gap in the familiar graph
guarantees a gap in the target graph. We use \emph{gap-preserving} as shorthand.

Each proof starts with a witness distribution \(\Pfamiliar\) on the familiar
graph and constructs a quantum distribution \(\Ptarget\) on the target graph.
A classical model for \(\Ptarget\) would give a classical model for
\(\Pfamiliar\), contradicting the choice of witness. The search can chain
several piggybacks until reaching a graph with an established gap, called a
\emph{seed}, then transfer the witness back along the chain.

\paragraph{Model conventions.}
The visible vertices of a graph \(G\) are laboratories. A directed arrow
between laboratories carries the sending laboratory's full classical output.
The directed arrows form an acyclic graph. Independent sources supply either
a shared random variable or a quantum state with one subsystem per recipient.
We call the set of recipients of a source its \emph{facet}. Each laboratory
operates on the systems and classical inputs available locally and reports an
outcome. We allow private randomness, finite output alphabets, and
finite-dimensional systems. A witness construction may enlarge an output
alphabet.

Write \(\Cset(G)\) for the distributions attainable with classical sources
and \(\Qset(G)\) for the distributions that also allow quantum states at
the designated quantum sources. For any choice of visible alphabets,
\(\Cset(G)\subseteq\Qset(G)\). A witness is a distribution in
\(\Qset(G)\setminus\Cset(G)\) for the chosen alphabets.

In a classical model, independent sources and one response kernel per
laboratory determine the joint distribution. The \emph{expanded DAG} draws
each source as a vertex with an arrow to each recipient. A response kernel
can also be represented by a deterministic response and an independent
private random seed. We draw a separate vertex for each private seed when
testing which inputs a predictor can learn.

Finite classical sources suffice for finite visible alphabets. Fix the laws
of all but one source and keep all response kernels fixed. Each value of the
remaining source gives a vector of observed probabilities. The observed
probability vector is an average of the conditional vectors, so
Carathéodory's theorem represents the average by a finite mixture. Applying
the same argument to each source in turn preserves the observed distribution.
With finitely many inputs, each response kernel can sample a finite table
with one output for each input configuration.

\paragraph{Reading the figures.}
Rounded squares mark visible laboratories. Red circles \(\lambda_i\) mark
classical sources, and green circles \(\psi_i\) mark quantum sources. Each
source arrow takes the colour of the source. Arrows between panels show
the direction of the graph edits in \eqref{eq:direction}.

A laboratory can prepare a source resource locally when the laboratory is
the only recipient. A source can also carry an independent classical register
for a subset of its recipients. A quantum source can carry an additional
independent state for a subset. We identify graphs that differ only by
local preparation or the choice to combine redundant sources. A figure may therefore draw a single source where
the proof uses separate resources.

\pageof{Reading the rules}\label{sec:map}
Each circle is a laboratory that reports a classical outcome. We call a laboratory that sends an arrow a \emph{parent} of its recipient, and the recipient its \emph{child}. An arrow from one circle to another allows the recipient to use that outcome. Each diamond is a new, independent source. A quantum source can send different parts of one state to several laboratories. A classical source can hand each recipient a copy of a random string. Each laboratory also has private randomness.
\duo{Quantum source and two shares}{\vtx{a}{0}{0}{a}\vtx{b}{3}{0}{b}\qsrc{q}{1.5}{1.1}{a,b}}
{Classical source and two copies}{\vtx{a}{0}{0}{a}\vtx{b}{3}{0}{b}\csrc{q}{1.5}{1.1}{a,b}}
A graph has a \emph{quantum--classical gap} if some observed distribution can be produced using its permitted quantum resources and cannot be produced when every source is classical. A graph with an independently established gap is a \emph{seed}. We consider finite experiments and may choose the size of each laboratory's output. A laboratory can therefore report a pair, including a setting or a teleportation correction in addition to its usual result. A claim about a graph establishes the existence of a gap for some choice of output sizes.

\begin{rulebox}{Find a rule by the change you want to make}
\begin{tabular}{@{}p{.47\linewidth}p{.45\linewidth}@{}}
Fix a laboratory's announcement & Opening triangle example, p.~1\\
Upgrade a source or add connections & \readat{sec:resources}{Compare resources}\\
Forget an announcement & \readat{sec:forget}{Relay or teleport}\\
Keep a selected outcome & \readat{sec:select}{Postselection}\\
Identify a setting with an outcome & \readat{sec:stitch}{Stitch two nodes}\\
Give one node two recorded outputs & \readat{sec:stitch}{Split a node}\\
Remove inputs by perfect prediction & \readat{sec:template}{The shared argument}\\
Use a parent as predictor & \readat{sec:parent}{Parent Fritz}\\
Use separation in the original graph & \readat{sec:separate}{Original Fritz}\\
Combine independence statements & \readat{sec:kpc}{The KPC example}\\
Use an entropy inequality & \readat{sec:entropy}{The same KPC example}\\
Delete edges using independence alone & \readat{sec:nopred}{Prediction-free deletion}
\end{tabular}
\end{rulebox}
\textbf{Why these proofs can be chained.} Let \(P'\) be an arbitrary witness for the simpler graph. Construct a distribution in the original graph from which a specified operation recovers \(P'\). Prove that applying that operation to a classical explanation would explain \(P'\) classically. The contradiction proves a gap for the original graph. The next rule can now use that gap even if it constructs a very different witness.

\trio{Forget a relay}{\vtx{a}{0}{0}{a}\vn[pre]{v}{1.7,0}{$v$}\vtx{b}{3.4}{0}{b}\de{a}{v}\de{v}{b}
\draw[newedge] (a) to[bend right=40] (b);}
{Select a joint outcome}{\vtx{a}{0}{0}{a}\vtx{b}{3.4}{0}{b}\vn[pre]{v}{1.7,1.5}{$v$}
\qsrc{l}{0}{1.5}{a,v}\qsrc{r}{3.4}{1.5}{v,b}}
{Predict a field}{\vtx{x}{0}{0}{X}\vtx{s}{3.4}{0}{S}\csrc{l}{1.7}{1.5}{x,s}
\node[lab] at (1.7,-.65) {\(S=g(X)\)};}
\captext{Three ways to recognize the next move: follow an arrow through a relay, look for two sources meeting at a measurement, or look for common data that two laboratories can both announce.}

\pageof{When the same experiment still works}\label{sec:resources}
\textbf{Upgrade a source.} Replace a shared classical source by a quantum one on the same recipients. One can also give a subset of its recipients an additional quantum source. A gap already present survives both changes.
\duo{Triangle with an upgraded source}{\vtx{a}{0}{0}{A}\vtx{b}{3.8}{0}{B}\vtx{r}{1.9}{2.8}{R}
\qsrc{ab}{1.9}{-.7}{a,b}\qsrc{u}{-.15}{1.75}{a,r}\csrc{v}{3.95}{1.75}{b,r}}
{Known triangle seed}{\trianglefig{0}}
\quant A quantum source can draw a classical label and send perfectly distinguishable copies. An optional extra source can send a fixed state. The old experiment still works.
\clas After treating all sources as classical, the new graph gives exactly the old possibilities. In the subset case the larger source may distribute the extra independent string too, with the other recipients ignoring it.

\textbf{Add connections within a classical equivalence class.} A gap passes to a graph with extra arrows or larger sources if both graphs allow the same classical distributions. If both quantum and classical sets agree, a gap passes in either direction.
\duo{The graph with an extra arrow}{\bellsettings{1}}{Bell settings with a common cause}{\bellsettings{0}}
\quant Ignore the extra arrow. Every old experiment still runs.
\clas Here \(X\) takes input only from the settings source and its private randomness. That source can sample \(X\) in advance and give it to both settings laboratories. Thus it can simulate the added arrow in a classical model. Equality of the classical sets ensures the old witness is still outside them.
\textbf{When can a source work equally well in either form?} Suppose all but at most one recipient have no visible parent and receive no other shared source. Each such laboratory makes a fixed measurement. The source can make those measurements itself, send the outcomes as classical data, and tell the remaining recipient which conditional state to prepare locally. This preserves the quantum statistics when the source becomes classical. A free setting that reads a single source in the Bell example has precisely this property \cite{Equiv}.

\begin{ideabox}{Why check the classical sets?}
Add every forward arrow in a time ordering of the laboratories of an ordinary Bell experiment. The enlarged classical graph can sample \emph{any} observed distribution, one outcome at a time conditional on all its predecessors. Even a Bell violation has a classical explanation with those extra communications.
\end{ideabox}

\pageof{Forget an announcement}\label{sec:forget}
Forgetting \(v\) means summing over every value \(v\) reported. A child can still receive information that \(v\) forwarded during the experiment. In this example, 0 sends its announcement to 1 and shares a separate source with 2.
\trio{Before forgetting 0}{\margfig{0}}{Ordinary relay}{\margfig{1}}{With teleportation}{\margfig{2}}
\captext{The source between 1 and 3 stays quantum in both outputs. Look at the new link between 1 and 2: the ordinary relay supplies shared classical data. Teleportation also permits a quantum state there. The visible edge \(1\to2\) stays in all three pictures.}

\textbf{Ordinary marginalization.} Give each visible child of \(v\) all of \(v\)'s visible parents. Every incident source keeps its other recipients and its type. Add a classical source for its other recipients together with all of \(v\)'s children. Add shared classical randomness over the children for \(v\)'s private random choice. A classical source contained in a larger classical or quantum source can be absorbed by it.
\quant Run the desired resources for the surviving recipients. Each old source can also send a random string to \(v\) and its other recipients. Let \(v\) report the strings, its visible inputs and a private seed. Every child can extract what it needs.
\clas A child with these inputs can recompute the old announcement of \(v\). Sampling its noise once for all children preserves their joint distribution.

\textbf{Marginalization with teleportation.} Let \(T\) be the children of \(v\) that already share a quantum source with \(v\). Each quantum source incident to \(v\) may now also distribute quantum shares to every member of \(T\).
\quant Place each destination subsystem at \(v\). For a child in \(T\), its old shared quantum source can additionally supply an entangled pair. Laboratory \(v\) measures its half together with the subsystem and sends the classical correction through the visible arrow. The child recovers its subsystem. Distinct destinations receive distinct subsystems.
\clas The ordinary classical explanation just constructed still fits the enlarged sources. In the picture, making the 1--3 source classical in the rightmost graph reaches an instrumental seed: 3 supplies a setting to 1 through shared data, 1 sends its outcome to 2, and 1 and 2 share a quantum state.

\pageof{Keep only selected runs}\label{sec:select}
Choose a node \(v\) and an outcome of positive probability. Keep runs with that outcome, then drop \(v\). Let \(B\) be its visible parents and let its \emph{siblings} be the laboratories sharing a source with it. Add a classical source for \(B\) and all siblings together, and a quantum source for all its quantum siblings. Existing sources keep their surviving recipients.
\duo{Laboratory 0 measures two shares}{\swapfig{0}}{Conditioned on one Bell outcome}{\swapfig{1}}
\captext{Here 0 has no visible parents. A joint measurement at 0 can leave 2 and 3 entangled on the selected runs. The other source, between 1 and 2, survives.}
\begin{rulebox}{Two checks before postselecting}
Every visible grandparent of \(v\) must also be its parent. Every source touching \(B\) must either include \(v\) or have all its recipients inside \(B\).
\end{rulebox}
\quant One source at \(v\) draws the desired shared label; the others guess it uniformly. Keep agreements. Every label has the same chance of acceptance, so its desired law survives. If there is no such source, \(v\) draws the label. Prepare the new quantum state at \(v\), teleport the outside shares, and select the outcomes needing no correction.

The parents announce uniform guesses. Laboratory \(v\) can run their intended measurements in causal order: their visible parents lie in \(B\), sources wholly inside \(B\) can be prepared locally, and every source reaching both \(B\) and the outside can deliver the parents' shares to \(v\). Keep runs where the guesses match the measured outputs. For each possible output record the chance of a match is the same. The accepted records and outside systems thus have the desired joint statistics.
\clas Condition the old sources feeding \(v\) or lying inside \(B\) on the selected outcome. The new classical source can draw this joint sample and the conditional parent outputs together. Give each parent its coordinate and each sibling its original source data. Other sources remain independent, and children of \(v\) can insert its selected value.
\textbf{A failed check.} Each parent \(p_i\) and its outside observer \(o_i\) report the same fair bit, independently for the two pairs. Select when the parents agree. The outside observers now share a fair bit. In the proposed graph the new source reaches only the parents, so \(o_1,o_2\) must be independent. The original sources each violate the second check.

\duo{Select when the parents agree}{
\vtx{o}{0}{1.6}{o_1}\vtx{p}{1.4}{0}{p_1}\vtx{v}{2.8}{-1.1}{v}
\vtx{r}{4.2}{0}{p_2}\vtx{t}{5.6}{1.6}{o_2}
\csrc{l}{1.4}{1.6}{o,p}\csrc{k}{4.2}{1.6}{r,t}\de{p}{v}\de{r}{v}}
{The attempted replacement}{
\vtx{o}{0}{1.6}{o_1}\vtx{p}{1.4}{0}{p_1}
\vtx{r}{4.2}{0}{p_2}\vtx{t}{5.6}{1.6}{o_2}
\csrc{l}{1.4}{1.6}{o,p}\csrc{k}{4.2}{1.6}{r,t}\csrc{u}{2.8}{-1.1}{p,r}}

\pageof{Stitch two nodes or split one}\label{sec:stitch}
\textbf{Stitching.} Choose \(x\) with no visible parents and no shared source. Choose \(a\ne x\) outside its descendants. Delete \(x\) and let its children receive \(a\) instead. This condition keeps the new arrows acyclic. The search chooses a childless \(a\).
\duo{Bell experiment with free setting \(x\)}{\stitchfig{0}}{Use the value \(a\) as that setting}{\stitchfig{1}}
\quant In the Bell experiment sample \(x\) independently and uniformly from the alphabet of \(a\). Run the desired right-hand strategy, feeding \(x\) to the former children wherever it asks for \(a\). On the diagonal \(x=a\) every instruction receives precisely the intended value. For each outcome \(a\) and each record \(w\) at the remaining laboratories, the distribution in the stitched graph is
\[
   P'(a,w)=\frac{P(x=a,a,w)}{P_X(a)}=|\mathcal A|\,P(x=a,a,w),
\]
where \(\mathcal A\) is the set of possible values of \(a\). The denominator is constant because \(x\) was chosen uniformly. Each run in the stitched experiment is represented on the diagonal with the same weight.
\clas Fix all classical random choices except \(x\). The value of \(a\) is already determined independently of \(x\). Wire that value into the former children of \(x\). Averaging these runs gives exactly the same normalized diagonal. A gap on the right therefore proves one on the left.

\textbf{Splitting.} Replace \(s\) by \(s,s'\). Each inherits its visible parents, its outgoing arrows, and membership in each of its sources. Also give the pair shared classical randomness, which can be carried by any source they already share.
\duo{One output can be a pair}{\splitfig{0}}{Two visible coordinates}{\splitfig{1}}
\quant One laboratory can carry out both split laboratories' measurements on their respective systems and announce the pair. Each child can read either coordinate.
\clas With one laboratory reporting a pair, its conditional distribution given its inputs can be sampled from one private random seed. Give the split laboratories the same seed and inputs. They compute the same pair and each reports its own coordinate. Children again receive both. Thus a gap for the split graph proves a gap for the original. Fritz uses this when one coordinate records a setting and the other records a measurement result.

\pageof{One construction, four reasons it works}\label{sec:template}
Pick a predictor \(X\) and a predicted output \(S\). We will engineer an experiment in which reading the \emph{whole announcement} at \(X\) tells us \(S\) perfectly. Which of the old inputs can a classical explanation dispense with? Consider this small example of the construction.
\trio{\(G\): choose \(X,S\)}{\siblingfig{0}}{\(G_1\): restrict \(S\)}{\siblingfig{1}}{\(G'\): forget \(X\)}{\siblingfig{2}}
\captext{In this slice of a larger graph, \(V\) sends an input that \(X\) cannot read. \(G_1\) deletes that arrow. The source shared by \(S,X\) supplies classical data to \(S\). With a childless \(X\) removed, a source serving \(S\) alone becomes private randomness.}
For a sibling predictor, \(S\) keeps its parents seen by \(X\): inputs that also enter \(X\), including shared sources. It loses every other parent. Each retained quantum source can still send a quantum state to its recipients other than \(S\), and sends shared classical data to the whole group. For a visible parent predictor use the simpler rule on the next page.
\quant Start with \emph{any} experiment witnessing a gap in \(G'\). Use the construction for forgetting a node to run it in \(G_1\). For a shared-source predictor, put \(S\)'s random seed in the shared classical data. Then \(X\) can compute \(S\) and announce \((x_{\rm old},s_{\rm prediction})\); its children still use \(x_{\rm old}\). A predicting parent can instead send the seed along its arrow, as on the next page. For several predictors, each reports the classical inputs it sees, so their combined record determines \(S\). This experiment also runs in \(G\), ignoring the deleted inputs.
\begin{rulebox}{What a hypothetical classical explanation must satisfy}
\textbf{P1.} The sources in the classical explanation are independent. Given them, the observed nodes are sampled in causal order, each from its own inputs in \(G\) and private randomness.\\
\textbf{P2.} \(S\) can be read without error from the full \(X\), written \(H(S\mid X)=0\).\\
\textbf{P3.} The observed announcements obey the observable independences of \(G_1\), with the predictor present. The constructed experiment can run there even with the extra field. Observable graphical separations hold for quantum experiments too \cite{HLP}.
\end{rulebox}
\clas We need a classical model of \(G_1\) for these same observations. The next four arguments give sufficient reasons it exists. Delete a childless predictor, or sum over a predictor with children and relay its inputs as on p.~\pageref{sec:forget}. A split predictor loses just its predicting copy. The result explains the original witness for \(G'\) classically, a contradiction. Only this final graph is handed to another piggyback.

\pageof{Fritz I. A parent predicts its child}\label{sec:parent}
\begin{rulebox}{If \(X\) sends its outcome to \(S\), prediction itself suffices}
Keep the arrow \(X\to S\). Delete every other input of \(S\), including its membership in sources shared with \(X\). The same holds for several predicting parents by keeping their arrows.
\end{rulebox}
\duo{A parent can also share a source}{\vtx{x}{0}{1.8}{X}\vtx{s}{3.4}{1.8}{S}\vtx{v}{3.4}{-.3}{V}
 \qsrc{l}{1.4}{.4}{x,s}\de{x}{s}\de[dropedge]{v}{s}}
{The restricted graph, before removal}{\vtx{x}{0}{1.8}{X}\vtx{s}{3.4}{1.8}{S}\vtx{v}{3.4}{-.3}{V}\de{x}{s}
 \node[lab,below] at (1.7,1.65) {send \(x\) and a seed};}
\quant In the restricted experiment \(S\) can sample its output as \(h(x,r)\) using a fresh random seed \(r\). Generate \(r\) at \(X\) and append it to \(x\). The child now outputs exactly the same \(h(x,r)\), and the predictor's full report determines it.
\clas P2 says \(S=g(X)\) on every event of positive probability. In any classical explanation replace \(S\)'s old instruction by this function. It uses a permitted input and preserves the full joint distribution. Here P1 and P2 suffice.

\begin{validbox}{Keep the quantum resource by predicting one field}
Take the chain \(0\to1\to2\to3\), with independent quantum sources on \(01,02,13\) (digits list recipients). Let 0 predict one field \(1'\) of the output at 1. Split 1 first. Before removing 0, the predicted field reads only 0's visible announcement. The original field 1 still receives its share of source \(13\). If we had restricted the entire laboratory 1, that share would have disappeared. The full sequence is drawn on p.~\pageref{sec:chain}.
\end{validbox}
After \(X\) leaves, its children may share classical data from the relay. This explains why shared sources can reappear in the final picture even though \(S\) kept only one arrow in the middle picture.

\pageof{Fritz II. What can independent inputs predict?}\label{sec:separate}
Write \(B\) for the inputs \(S\) keeps and \(A\) for the inputs it loses, including its own private noise. In the original graph, expose each hidden source as a node. Check whether the predictors outside \(B\) are separated from \(A\) once \(B\) is known. The usual path check is called \emph{d-separation}: a path closes at a conditioned ordinary intermediate node, or at a meeting of two arrowheads unless that meeting or a descendant is conditioned on.
\begin{rulebox}{The classical direction in a four-cell picture}
Fix the kept inputs at \(B=b\). Classical instructions give \(S=h(b,A)\), and P2 gives \(S=g(X)\). Graphical separation makes \(A\) and \(X\) independent at this fixed \(b\). For two possible values of each, the common entry would have to satisfy
\[
\begin{array}{c|cc}
 & X=x_0 & X=x_1\\ \hline
 A=a_0 & h(b,a_0)=g(x_0)&h(b,a_0)=g(x_1)\\
 A=a_1 & h(b,a_1)=g(x_0)&h(b,a_1)=g(x_1)
\end{array}
\]
Independence puts every combination of positive-probability row and column on the table. All entries agree. Thus \(S\) can use a function of \(b\) alone. Apply this for each \(b\). The same reasoning applies to larger alphabets, almost surely.
\end{rulebox}
\quant The shared input gives the predictor the data needed to announce \(S\) as on p.~\pageref{sec:template}. The separation check is a fact about the original graph used for the classical argument. That proof needs no observable independence from P3. With several targets, replace one instruction at a time. Each replacement preserves the joint distribution, so the next prediction still holds.
\textbf{Example: the tetrahedron.} There is an independent source for each triple of laboratories and no visible arrows. Laboratory 0 predicts 1. Read each row as one source, each mark as one recipient.
\duo{Original sources}{\cols{0/0,1/1,2/2,3/3}
\members{-.35}{\(012\)}{0,1,2}\members{-1.0}{\(013\)}{0,1,3}
\members{-1.65}{\(023\)}{0,2,3}\members{-2.3}{\(123\)}{1,2,3}}
{After restriction and removing 0}{\cols{1/1,2/2,3/3}
\classmembers{-.35}{classical}{1,2}\classmembers{-1.0}{classical}{1,3}\members{-1.65}{quantum}{2,3}}
\captext{Fix the sources \(012,013\) seen by both 0 and 1. The unseen source \(123\) is independent of the remaining inputs at 0. The new classical links carry the shared data, and 2 and 3 retain a quantum source. This is the triangle seed.}

\pageof{The triangle hides two Bell settings}\label{sec:triangle}
Fritz's original experiment \cite{Fritz} illustrates why predicting two outputs together matters. Let \(A,B\) share an entangled state. A second source supplies a random setting \(x\) to \(A,R\); a third independently supplies \(y\) to \(B,R\). Ask the three laboratories to announce \((x,a)\), \((y,b)\), and \((x,y)\) respectively. Here \(a,b\) are the two measurement outcomes.
\duo{The three laboratories}{\trianglefig{0}}
{Show each setting as its own output}{\trianglefig{1}}
\captext{The right panel shows the restricted experiment. The small laboratories \(x,y\) are the setting fields split from \(A,B\). The record at \(R\) predicts both. The two blue sources are explicitly classical; an all-quantum triangle can run this experiment too.}
\duo{Forget the record at \(R\)}{\trianglefig{2}}
{A Bell test on the surviving fields}{
\node[draw=ccol,rounded corners,fill=ccol!6,align=center] at (2,2.45) {independent choices \(x,y\)};
\node[draw=qcol,rounded corners,fill=qcol!7,align=center] at (2,.95) {measure the shared state};
\node[align=center] at (2,-.15) {read \(a,b\) at settings \(x,y\)};}
\quant Choose any Bell experiment in the final graph. The two settings sources can include the random seeds generating the recorded settings. In the triangle \(R\) receives these seeds and reports both settings. Forgetting \(R\) leaves exactly the Bell distribution.
\clas In a classical triangle, the field \(x\) might initially use the hidden source between \(A\) and \(B\). Given the shared input of \(A,R\), all remaining inputs of \(A\) are independent of the remaining inputs at \(R\). If both always report the same \(x\), the table on p.~\pageref{sec:separate} forces \(x\) to depend only on that shared input. Apply the same argument to \(y\) on the other side. The two settings are generated independently of one another and of the source between \(A,B\). A classical explanation would now give a Bell-local model for the surviving fields, impossible for a Bell-violating experiment.
\begin{validbox}{A gap you can read directly from the triangle records}
Choose independent fair setting bits \(x,y\). Count a win when \(a\oplus b=xy\), where \(\oplus\) is addition modulo two. A classical Bell strategy wins at most three of the four setting pairs on average: for fixed local responses, adding the four left sides modulo two gives 0 (each response appears twice), while adding the four right sides gives 1. Shared random choices only average these bounds. On the state \((|00\rangle+|11\rangle)/\sqrt2\), measuring the Pauli observables \(Z\) or \(X\) at \(A\) and \((Z+X)/\sqrt2\) or \((Z-X)/\sqrt2\) at \(B\) wins with probability \((2+\sqrt2)/4>3/4\). The eigenvalues \(+1,-1\) are recorded as \(0,1\). Laboratory \(R\)'s two predictions certify the setting fields used to score the game.
\end{validbox}

\pageof{Fritz III. What the extra independence buys}\label{sec:kpc}
Khanna, Pusey and Colbeck (KPC) \cite{KPC} give a graph where the last path test fails. Use their names: \(C,D,E,F\) are observed. Source \(A\) feeds \(E,F\) and source \(B\) feeds \(D,F\). Ask \(F\) to predict \(E\). The restricted \(E\) keeps only the source shared with \(F\). Retain the ordinary output \(F\) by forgetting just its predicting field.
\trio{Original graph \(G\)}{\kpcfig{0}}{Restricted graph \(G_1\)}{\kpcfig{1}}{Classical explanation}{\kpcfig{2}}
\captext{Left: fixing \(A\) leaves \(F\leftarrow B\to D\) open. Middle: removing the two dashed arrows gives \(E\ind(C,D)\), while \(D,F\) may still measure a shared quantum state. Right: if the original observations had a classical explanation, the proof below would give this model with independent classical sources \(A',B\).}
\quant Begin with that Bell experiment: \(C\) supplies a setting to \(D\), the shared classical source supplies the setting \(E\) to \(F\), and \(D,F\) share a quantum state. Give \(F\) a copy of the data determining \(E\), including its random seed. In the original graph the extra arrows into \(E\) can be ignored. This realizes P2 and the observable independence in P3. In the deductions below \(F\) denotes the full report including its extra field.
\clas Write \(U=(C,D)\). In a proposed classical explanation, \(J\ind K\mid L\) means that once \(L\) is known, learning \(J\) gives no further information about \(K\).
\begin{enumerate}
\item \(E\ind B\mid U\) follows in \(G\): given \(C,D\), \(E\) uses only the independent source \(A\) and its own noise. P3 gives \(E\ind U\). Combining these yields \(E\ind(B,U)\).
\item \(F\ind U\mid B\) follows in \(G\): fixing \(B\) separates the inputs generating \(F\) from \(C,D\). P2 gives \(E\ind U\mid B,F\). Combine the statements to get \((F,E)\ind U\mid B\), then condition on \(E\) to get \(F\ind U\mid B,E\).
\item \(B\ind C\) in \(G\). The two preceding lines now give the factorization
\[
p(b,c,d,e,f)=p(b)\,p(c)\,p(d\mid b,c)\,p(e)\,p(f\mid b,e).
\]
\end{enumerate}
Sample \(A'\) independently with distribution \(p(e)\), let \(E=A'\), and give \(F\) the inputs \(A',B\). The factorization proves these instructions preserve all observed statistics. Forget \(F\)'s extra field. This is the desired classical explanation of the Bell experiment and completes the contradiction.

\pageof{Why the old source can be the wrong one}\label{sec:fresh}
Let \(A,B,C\) be independent fair bits. The symbol \(\oplus\) means addition modulo two. Keep the KPC arrows and give the laboratories the following instructions:
\[
 D=C\oplus B,\qquad E=A\oplus C\oplus D=A\oplus B,\qquad F=(A,B).
\]
Laboratory \(E\) can use the first expression because it receives \(A,C,D\). Laboratory \(F\) can predict it by the second. For each fixed \(C,D\), the fresh fair bit \(A\) makes \(E\) fair, so \(E\ind(C,D)\). Could \(E\) simply toss a private coin after reading the old \(A\)? Given \(A\), \(E\) and \(B\) are perfectly correlated or anticorrelated: \(E=A\oplus B\). A rule using only \(A\) and independent private noise would give \(E\ind B\mid A\). Thus it cannot preserve this classical joint distribution. The table constructs a different joint distribution with the same observations.
\duo{Sample \(A,B,C\) independently}{
\node[draw=qcol,fill=qcol!6,rounded corners,align=left] at (1,1) {\(\begin{array}{cc|cc} A&B&E&F\\\hline 0&0&0&00\\0&1&1&01\\1&0&1&10\\1&1&0&11\end{array}\)};}
{Sample \(A',B,C\) independently}{
\node[draw=ccol,fill=ccol!6,rounded corners,align=left] at (1,1) {\(\begin{array}{cc|cc} A'&B&E&F\\\hline 0&0&0&00\\1&1&1&01\\1&0&1&10\\0&1&0&11\end{array}\)};}
\captext{Each row has probability \(1/4\) and is repeated for \(C=0,1\) with equal weights. Corresponding rows have identical \(B,E,F\), hence identical \(C,D,E,F\). On the right use \(E=A'\), \(D=C\oplus B\), and \(F=(A'\oplus B,B)\).}
\begin{rulebox}{Two sufficient ways to prove a restricted graph has a classical explanation}
\textbf{Keep the old sources.} Use P1, P2 and P3 to deduce the local independence statements of the restricted graph. They say each node can be sampled given its permitted parents. This test can certify restrictions directly.

\textbf{Replace one source.} When one predicted node \(S\) keeps exactly one parent, a hidden source \(L\), remove \(L\) on paper. Make \(S\) a root and send its value to every other former recipient of \(L\). Prove the local independence statements of this graph over the remaining variables. Then a fresh independent source carrying \(S\) realizes the restricted graph, exactly as on p.~\pageref{sec:kpc}.
\end{rulebox}
\textbf{How the independence test works.} Start with graphical independences of the original classical model, the observable ones from the restriction, and the fact that once the predictor is given, learning other variables adds no information about \(S\). This last fact continues to hold with extra variables conditioned on. Close the list under symmetry, dropping variables from one side, moving part of that side into the conditioning set, and the reverse combining step (contraction). These are the \emph{semigraphoid rules}. The proof on the previous page is a short instance. This calculation uses a consequence of perfect prediction; independence from other variables given \(X\) alone also holds for an unrelated fresh random bit.

\pageof{Fritz IV. Read the same proof in bits}\label{sec:entropy}
Read \(H(E\mid F)=0\) as zero uncertainty about \(E\) once \(F\) is known. Read \(I(J;K\mid L)=0\) as independence of \(J,K\) given \(L\). Both quantities are nonnegative for the joint distribution of a classical explanation.
\begin{rulebox}{An entropy calculation for the KPC graph}
Again write \(U=(C,D)\). P1 supplies \(I(E;B\mid U)=0\) and \(I(F;U\mid B)=0\). P3 supplies \(I(E;U)=0\). The chain rule gives
\[
 I(E;B,U)= I(E;U)+I(E;B\mid U)=0.
\]
P2 supplies \(H(E\mid F)=0\). Consequently
\[
 0\le I(E;U\mid B,F)\le H(E\mid B,F)\le H(E\mid F)=0.
\]
Use the chain rule in its two orders:
\[
 \underbrace{I(F;U\mid B)+I(E;U\mid B,F)}_{0}
     = I(E;U\mid B)+I(F;U\mid B,E).
\]
Both terms on the right must vanish. These are the same ingredients that allow the fresh-source factorization on p.~\pageref{sec:kpc}.
\end{rulebox}
\quant The experiment on p.~\pageref{sec:kpc} again lets \(F\) announce \(E\) exactly.
\clas For the general test, take one entropy coordinate for every nonempty subset of the classical variables. Impose the Shannon inequalities (nonnegative conditional entropies and conditional mutual informations), the original graph's local independence equalities, full \(H(S\mid X)=0\), and the visible independence equalities of \(G_1\). A local equality for node \(v\) says \(I(v;N\mid\Pa(v))=0\), where \(N\) contains its non-descendants other than itself and its parents.
\textbf{Why one linear program suffices for a proposed conclusion.} Sum the local quantities for the chosen graph, calling the sum \(T\). Each summand is nonnegative. To test if all vanish, ask whether the constraints permit \(T\ge1\). All constraints before the last are homogeneous, so any positive \(T\) could be scaled to 1. Infeasibility forces \(T=0\) for every compatible classical model of \(G\). Each required independence then holds, so the construction on p.~\pageref{sec:fresh} gives a classical explanation in \(G_1\). The solver can certify \(T\le0\) by combining the premise equalities and nonnegative multiples of the inequalities. One solve checks each proposed set of conclusions; testing both choices on p.~\pageref{sec:fresh} can require two. Feasibility leaves the question undecided because a Shannon vector need not come from a distribution.

\begin{ideabox}{Why agreement on the census does not make the two tests equivalent}
For four arbitrary variables, expand each information quantity into entropies. The following two sums agree term for term after cancellation:
\begin{align*}
 &I(A;B\mid C,D)+I(C;D\mid A)+I(C;D\mid B)+I(A;B)\\
={}&I(C;D\mid A,B)+I(A;B\mid C)+I(A;B\mid D)+I(C;D).
\end{align*}
If the first four terms vanish, the last four vanish too. The semigraphoid rules applied to those premises derive none of the four conclusions \cite{Repo}.
\end{ideabox}

\pageof{Which piece of the announcement is predicted?}\label{sec:modes}
A laboratory that records \((s_{\rm measurement},s_{\rm setting})\) can have its second field predicted while its first is still a quantum measurement outcome. The predictor can likewise keep a full result when only a field of it is needed. The two choices are independent.
\begin{center}
\begin{tabular}{@{}p{.20\linewidth}p{.34\linewidth}p{.36\linewidth}@{}}
\toprule & \textbf{Predict the whole target} & \textbf{Predict a target field}\\\midrule
Remove the predictor & Tetrahedron: 0 leaves after predicting 1. & Triangle: \(R\) leaves after predicting the two setting fields.\\[5pt]
Keep the predictor & KPC: \(F\) stays available as a Bell outcome. & Split both: a predicting copy leaves and a target copy carries the setting.\\\bottomrule
\end{tabular}
\end{center}
\textbf{Keeping a predictor that has children.} Make a copy with the same parents, sources and children. Restrict the target using the copy as predictor. Sum over the copy so that its children still have access to whatever it could have relayed. For a childless predictor the predicting copy can simply be deleted. The original laboratory stays, with its ordinary output. When reading P2, \(X\) denotes the full predicting announcement available during the proof.
\duo{Split a predictor with a child}{\vtx{x}{0}{1.6}{X}\vtx{xp}{3.3}{1.6}{X'}\vtx{c}{1.65}{-.5}{c}
\qsrc{l}{1.65}{2.65}{x,xp}\de{x}{c}\de{xp}{c}}
{After summing over the copy}{\vtx{x}{0}{1.6}{X}\vtx{c}{1.65}{-.5}{c}
\csrc{l}{2.4}{1.85}{x,c}\de{x}{c}}
\captext{A local view of the relay with no teleportation channel to \(c\). The source shared by the copies can now supply classical data to \(c\). Other recipients of that source retain their shares.}
\textbf{Predicting a target field.} First split the target into two laboratories with the same inputs. Restrict only the copy. In the single-laboratory experiment the original device records the two outputs together. The predictor can report the setting while the original field records the measurement result.
\begin{ideabox}{Several targets and several predictors}
One predictor can announce several predictions in one step, as \(R\) does in the triangle. For the graphical test each target must pass its own check in the graph before any restrictions. A set of predictors can also see more than any one member: their combined report can determine \(S\). Each member must have a channel to \(S\) in the construction used here. The four-node census tries single predictors. Its stronger independence and entropy checks are applied to single targets.
\end{ideabox}

\pageof{A copy predicts another copy}\label{sec:twocopies}
Start with \(0\to1\to2\) and three quantum sources, on \(02,03,13\). Laboratory 0 has a child. We want 0 to keep its ordinary output, and we want to predict only one field of 3. Give 0 a copy \(0'\) with the same child 1. Give 3 a copy \(3'\). Now \(0'\) can announce the setting field \(3'\) while both original laboratories remain available.
\duo{Inputs of the copies before restriction}{
\cols{0/0,1/{0'},2/1,3/2,4/3,5/{3'}}
\members{-.4}{\(02\)}{0,1,3}\members{-1.2}{\(03\)}{0,1,4,5}\members{-2}{\(13\)}{2,4,5}}
{Inputs after restricting \(3'\)}{
\cols{0/0,1/{0'},2/1,3/2,4/3,5/{3'}}
\members{-.4}{\(02\)}{0,1,3}\members{-1.15}{\(03\)}{0,1,4}\classmembers{-1.9}{\(03\)}{0,1,4,5}
\members{-2.6}{\(13\)}{2,4}}
\captext{The row names refer to the sources in the original graph. Filled dots carry quantum systems and open squares carry classical data. In both panels the arrows are \(0\to1\), \(0'\to1\), \(1\to2\). The two copies share source \(03\). Source \(13\) still serves the original 3 after \(3'\) loses it.}

\textbf{Why the prediction licenses the deletion.} Condition on the shared input \(03\). The source \(13\) can reach the predictor's child 1, but the path \(0'\to1\leftarrow13\) meets two arrowheads there and we have conditioned only on the shared source. The other paths from \(13\) to \(0'\) also close at a meeting of arrowheads or at the conditioned source. The private noise of \(3'\) has no route back to \(0'\). The table proof on p.~\pageref{sec:separate} applies.

\textbf{Where the copy's other inputs go.} Sum over \(0'\). It shares no quantum source with its child 1, so it can relay its inputs to 1 classically. The result has arrows \(0\to1\to2\), quantum sources on \(02,03,13\), and classical sources on \(012\) and \(0133'\). In particular the larger classical source can carry to 1 the data used in predicting \(3'\).
\begin{validbox}{Finish the proof}
Postselect on 3, which has no visible parents. Its shares of the \(03\) and \(13\) sources give 0 and 1 a quantum link. Now 0 can teleport its share of \(02\) to its child 1. Forgetting 0 leaves \(1\to2\), a quantum source on \(12\), and a classical setting shared by \(1,3'\): an instrumental seed. Changing the order loses this route, since 0 and 1 had no quantum link before the postselection.
\end{validbox}
The copy \(0'\) needs to determine \(3'\), but can carry other information to 1. Deleting those other inputs to \(0'\) would require a second argument. This is the operational reason for using the one-way condition \(H(3'\mid 0')=0\).


\pageof{Deletion using independence alone}\label{sec:nopred}
KPC's prediction-free rule asks whether the extra visible independences of an edge-deleted graph are already enough. There is no predictor to remove. Draw the graph \(H\) obtained by deleting the proposed arrows from \(G\); keep the same source model. Test whether the original classical graphical independences, together with the observable independences of \(H\), imply the local classical independence requirements of \(H\).
\duo{Direct and indirect routes}{\vtx{x}{0}{0}{X}\vtx{y}{1.8}{1.7}{Y}\vtx{z}{3.6}{0}{Z}
\de{x}{y}\de{y}{z}\de[dropedge]{x}{z}}
{Only the route through \(Y\)}{\vtx{x}{0}{0}{X}\vtx{y}{1.8}{1.7}{Y}\vtx{z}{3.6}{0}{Z}\de{x}{y}\de{y}{z}}
\captext{A minimal example of the test. In the graph on the right the observable condition is \(Z\ind X\mid Y\). These three observed variables suffice to display the deduction; a larger graph may also have hidden sources.}
\quant An experiment on \(H\) runs on \(G\) by ignoring the deleted arrows. Its visible distribution obeys every observable separation of \(H\).
\clas In this example any classical distribution on the left factors as \(p(x)p(y\mid x)p(z\mid x,y)\). The extra observable condition allows the last factor to be written \(p(z\mid y)\) on its support. That is a classical model on the right. In general, derive all the needed local independences of \(H\) by the rules on p.~\pageref{sec:fresh}; their joint validity gives the corresponding factorization, even when classical hidden variables occur. Whenever \(H\) has a gap, this proof transfers it to \(G\). The simple three-circle example illustrates the implication between models; a gap is a separate premise.

\begin{rulebox}{Try the same idea on the KPC example}
Retain \(E\ind(C,D)\) but remove the prediction of \(E\) by \(F\). On p.~\pageref{sec:kpc} the second deduction needed \(E\ind U\mid B,F\). Observing \(E\ind U\) gives no such conditional statement. The repository checks this prediction-free version of the KPC deletion and the test fails. In contrast, the three-circle example above needs just the one observable independence.
\end{rulebox}
The code can run this test with an empty prediction list. The default census uses the four Fritz variants described above.

\pageof{Follow a complete proof to an instrumental seed}\label{sec:chain}
Start with the chain and sources from the parent example. Below, marks in one row are recipients of one independent source. Filled circles denote quantum shares and open squares denote copies of classical data. The visible arrows are given under each panel.
\duo{1. Start with \(0\to1\to2\to3\)}{\cols{0/0,1/1,2/2,3/3}
\members{-.4}{\(01\)}{0,1}\members{-1.2}{\(02\)}{0,2}\members{-2}{\(13\)}{1,3}}
{2. Split 1; restrict \(1'\) to parent 0}{\cols{0/0,1/1,2/{1'},3/2,4/3}
\members{-.4}{\(01\)}{0,1}\members{-1.2}{\(02\)}{0,3}\members{-2}{\(13\)}{1,4}}
\captext{Panel 2 keeps \(0\to1\to2\to3\) and adds \(0\to1'\to2\). Only the predicted field has lost its source inputs. Its perfect prediction is an intermediate requirement.}
\duo{3. Forget the predictor 0}{\cols{0/1,1/{1'},2/2,3/3}
\members{-.4}{quantum}{0,2}\members{-1.2}{quantum}{0,3}\classmembers{-2}{classical}{0,1,2}}
{4. Forget 1 using teleportation}{\cols{0/{1'},1/2,2/3}
\members{-.4}{quantum}{1,2}\classmembers{-1.2}{classical}{0,1}}
\captext{Panel 3 has \(1\to2\to3\) and \(1'\to2\). Panel 4 has \(1'\to2\to3\). The original field 1 shared a source with its child 2, so it could teleport to 2 its share of the state with 3.}
\textbf{Read the panels backwards to prove the gap.} Panel 4 is an instrumental seed. Its setting \(1'\) can be sampled independently, the treatment 2 receives it and sends its outcome to 3, and 2 and 3 can measure a shared quantum state. The classical source over \(1',2\) could sample the setting and any randomness associated with it. This is the instrumental variant with both a setting edge and a shared setting source in the seed list \cite{Repo}.

Every witness for panel 4 can therefore be run in panel 3 using teleportation. A classical explanation in panel 3 would give one in panel 4 by forgetting 1. The parent Fritz argument then carries the gap from panel 3 to panel 1: split the target, arrange its prediction, and forget the predictor. Panel 2 carries the temporary prediction. The gap proved for panel 1 can be used in a new argument with any witness.

\pageof{What the search actually checked}\label{sec:record}
The four-node search starts with 426 candidate classes after excluding the 20 classes whose distributions already have classical explanations on a graph with no hidden source. The table gives the ten stages of the depth-one run (\texttt{--depth-one}). The deeper default search finds more intermediate graphs and proves the same 380 input classes \cite{Repo}. A class here groups graphs connected by proven equivalences preserving both quantum and classical observed distributions. The elementary rules and known starting cases establish 374 classes; 52 go on to the Fritz stages.
\begin{center}
\begin{tabular}{@{}lrr@{}}\toprule
First successful justification & Newly established classes & Cumulative\\\midrule
Parent predictor & 3 & 377\\
Graphical separation & 2 & 379\\
Conditional-independence deductions & 1 & 380\\
Entropy test on the same candidates & 0 additional & 380\\\bottomrule
\end{tabular}
\end{center}
\textbf{Classes and candidate steps answer different questions.} At this search depth, 696 candidates reached the comparison after failing the graphical test. The two stronger calculations both accepted the same 37 and left 659 uncertified. Several successful steps may concern one class or lead to graphs whose gaps are still open. The equality of the 37 decisions is a measured result for this set. The entropy test can use all Shannon inequalities and full determinism; the independence calculation uses a smaller collection of deductions.
\textbf{How a candidate is chosen.} For each predictor, first consider its children in the parent stage. The other stages consider nodes sharing a source with the predictor while excluding its children. A predictor that is both a parent and a sibling belongs to the first stage. The default also skips descendants of the predicted node: its private noise reaches any such predictor by a directed path, defeating the direct separation test. A descendant with a shared channel can still be certified by a stronger test. For each pair the kept inputs determine one deletion. The default stronger search tries the fresh-source conclusion; the direct, original-source conclusion is available as an option.
\begin{center}\small
\begin{tabular}{@{}>{\raggedright\arraybackslash}p{.40\linewidth}>{\raggedright\arraybackslash}p{.52\linewidth}@{}}\toprule
Where to check the argument & Repository entry point\\\midrule
Elementary transformations; predictor pool, deletion, test and removal & \texttt{quantum\_mDAG.py}: \texttt{fritz\_pool}, \texttt{fritz\_deletion}, \texttt{fritz\_certificate}, \texttt{fritz\_realise}\\
Independent deductions; summed entropy test & \texttt{semigraphoid.py}; \texttt{entropic\_lp.py}, \texttt{EntropicLP.implies\_all}\\
Four-node stages; KPC comparison & \texttt{tests/test\_baseline\_slow.py}, \texttt{tests/test\_fritz\_entropic.py}, \texttt{tests/test\_kpc\_figure4.py}\\\bottomrule
\end{tabular}
\end{center}
\textbf{Reading the rules in reverse.} Start from a seed and look for graphs whose forward transformation gives that seed. For example, insert a relay on a visible edge or attach a laboratory which can announce a constant. Checking the forward conditions certifies each proposed predecessor. This enumerates more graphs with a gap; it uses the same implication proved in the forward presentation.

\textbf{Reproduce the illustrated examples.} In the directory containing this PDF, run \texttt{python check\_examples.py} with the repository dependencies installed. It checks the labelled transformations, both complete chains, the KPC deductions and the two binary distributions. Run \texttt{./build.sh} to rebuild the PDF.

\clearpage\markboth{Sources}{}
\begin{thebibliography}{9}
\bibitem{Fritz} T. Fritz, \emph{Beyond Bell's theorem: correlation scenarios}, New Journal of Physics 14, 103001 (2012). The triangle construction embeds Bell settings in recorded outcomes. \href{https://doi.org/10.1088/1367-2630/14/10/103001}{Article}.
\bibitem{HLP} J. Henson, R. Lal and M. F. Pusey, \emph{Theory-independent limits on correlations from generalized Bayesian networks}, New Journal of Physics 16, 113043 (2014). Observable graphical separation in quantum networks. \href{https://doi.org/10.1088/1367-2630/16/11/113043}{Article}.
\bibitem{Evans} R. J. Evans, \emph{Graphs for margins of Bayesian networks}, Scandinavian Journal of Statistics 43, 625--648 (2016). The classical graph after a node is marginalized.
\bibitem{KPC} R. Khanna, M. F. Pusey and R. Colbeck, draft manuscript and accompanying notebook, Lemma 1 and Corollary 3. The repository reproduces the six-variable example, the independence argument, and the comparison with the draft's Figure 4. \href{https://github.com/eliewolfe/mDAG-analysis/blob/6af05cb/manuscript/kpc_figure4_vs_code.md}{Worked comparison}.
\bibitem{Repo} \emph{Piggybacks for quantum-classical gaps}, mDAG-analysis, main revision \texttt{6af05cb}. Sections 0 and 2--5 give the elementary rules; Sections 6--8 give the Fritz constructions and their certificates; Appendix B gives the entropy test. \href{https://github.com/eliewolfe/mDAG-analysis/blob/6af05cb/manuscript/piggybacks.md}{Technical manuscript}.
\bibitem{Equiv} \emph{Observational equivalence rules for quantum causal models}, companion illustrated note. Section 5 gives examples of transferring a gap within an equivalence class. \href{https://github.com/eliewolfe/mDAG-analysis/blob/6af05cb/manuscript/observational\%20classification/observational_classification.pdf}{Illustrated PDF}.
\end{thebibliography}

\end{document}
